% Source: Gerzensee "Calculus" slides 43-54 (the representative firm and the OLS
% problems; the coordinate-by-coordinate idea; partial derivatives and
% Definition 3.5; the Cobb-Douglas example; the quadratic-form gradient
% (A+A^T)x of Exercise 3.3/Example 11.2; the tangent-plane picture and the total
% differential; the first-order condition and its concave sufficiency; the four
% surfaces x^2+y^2, -x^2-y^2, x^2-y^2, -x^3; the representative firm's FOCs and
% factor shares).
% Supplementary: Fukuda (2026), Ch. 11 -- Definition 11.1 (differentiability),
% Prop. 11.3 (differentiability implies continuity), Remarks 11.3-11.5
% (counterexamples), Prop. 11.4 (C^1 implies differentiable), Thm 11.1 (chain
% rule), Prop. 11.6 (gradients and level sets), Def. 11.3 (homogeneity),
% Thm 11.2 (Fermat), Thm 11.3 (differentiable convexity), Def. 11.6 (Hessian),
% Thm 11.5 (second-order Taylor), Thm 11.6 (Hessian and convexity),
% Thm 11.7 (implicit function theorem, p. 419).
% External references actually cited in this chapter: Rudin (1976).

\chapter{Differential Calculus in Several Variables}
\label{ch:calculusn}

\section{Motivation: Two Multivariate Problems}\label{sec:cn-motivation}

Two problems drive the chapter \autocite[slide 43]{fukuda2026calculus}. The
representative firm chooses capital and labour to solve
\begin{equation}\label{eq:rep-firm}
	\max_{K,L\geq0}\ \Pi(K,L),\qquad
	\Pi(K,L)\eqdef AK^{\alpha}L^{1-\alpha}-(rK+wL),
\end{equation}
and the econometrician fits a line by
\begin{equation}\label{eq:ols-problem}
	\min_{(\beta_1,\beta_2)\in\R^2}\ \sum_{i=1}^{n}\bigl(y_i-\beta_1-\beta_2x_i\bigr)^2 .
\end{equation}

Neither is harder than the one-variable problems of \autoref{ch:calculus1} once
the right notion of derivative is in place, and finding that notion is the
substance of the chapter. Two distinct ideas are involved and are frequently
conflated. The first is to vary one coordinate at a time, which produces partial
derivatives. The second is to approximate $f$ near a point by an affine function
of the whole vector, which produces the total derivative. In one variable the two
coincide. In several variables the second is strictly stronger, and it is the
second that supports the chain rule, the first-order condition, and the implicit
function theorem.

\section{Partial Derivatives and the Gradient}\label{sec:cn-partial}

\begin{definition}[Partial derivative and gradient]\label{def:partial}
	Let $f\colon A\to\R$ with $A\subseteq\R^n$ and let $x$ be interior to $A$. For
	$i\in\{1,\dots,n\}$ the \dfnas{$i$th partial derivative}{partial derivative} of
	$f$ at $x$ is
	\[
		\frac{\partial f}{\partial x_i}(x)\eqdef
		\lim_{h\to0}\frac{f(x_1,\dots,x_i+h,\dots,x_n)-f(x_1,\dots,x_i,\dots,x_n)}{h},
	\]
	provided the limit exists and is finite; equivalently, it is the ordinary
	derivative at $x_i$ of the section $t\mapsto f(x_1,\dots,t,\dots,x_n)$. When all
	$n$ partial derivatives exist, the \dfn{gradient} of $f$ at $x$ is the vector
	\[
		\nabla f(x)\eqdef
		\left(\frac{\partial f}{\partial x_1}(x),\dots,
		\frac{\partial f}{\partial x_n}(x)\right)^{\transp}\in\R^n .
	\]
	\nidx{partial}{$\partial f/\partial x_i$, partial derivative}
	\nidx{nabla}{$\nabla f$, gradient}
\end{definition}

This is Definition~3.5 of \textcite{fukuda2026notes}
\autocite[slides 46, 48]{fukuda2026calculus}, where the alternatives
$\partial f(x)/\partial x_i$, $f_{x_i}(x)$ and $f'(x)$ for the gradient are also
recorded. This volume writes the gradient as a column vector, consistently with
the front-matter chapter on notation and conventions; the row vector
$\nabla f(x)^{\transp}$ is the Jacobian of $f$ viewed as a map into $\R^1$, and
conflating the two is the most frequent source of transpose errors in
comparative statics.

\begin{eg}[Cobb--Douglas]\label{eg:cobb-douglas-partials}
	For $f(x,y)=x^{\alpha}y^{1-\alpha}$ with $\alpha\in(0,1)$ and $x,y>0$, fixing
	$y$ makes $f$ a power function of $x$, so
	\[
		\frac{\partial f}{\partial x}=\alpha x^{\alpha-1}y^{1-\alpha}
		=\alpha\left(\frac{y}{x}\right)^{1-\alpha},
		\qquad
		\frac{\partial f}{\partial y}=(1-\alpha)x^{\alpha}y^{-\alpha}
		=(1-\alpha)\left(\frac{x}{y}\right)^{\alpha}.
	\]
	Both are positive and each is decreasing in its own argument: marginal products
	are positive and diminishing. Both are homogeneous of degree zero, so they depend
	on the two inputs only through the ratio $x/y$ --- the property that makes the
	factor-price equations of \autoref{sec:cn-firm} determine the capital--labour
	ratio and nothing else \autocite[slide 47]{fukuda2026calculus}.
\end{eg}

\begin{eg}[Gradient of a quadratic form]\label{eg:quadratic-gradient}
	Let $A=(a_{ij})\in M_{n\times n}(\R)$ and $f(x)=x^{\transp}Ax
		=\sum_{i,j}a_{ij}x_ix_j$. Differentiating with respect to $x_k$, the terms
	containing $x_k$ are $a_{kk}x_k^2$ and, for $j\neq k$, $a_{kj}x_kx_j$ and
	$a_{jk}x_jx_k$, so
	\[
		\frac{\partial f}{\partial x_k}(x)=2a_{kk}x_k+\sum_{j\neq k}(a_{kj}+a_{jk})x_j
		=\sum_{j=1}^{n}(a_{kj}+a_{jk})x_j ,
	\]
	that is,
	\begin{equation}\label{eq:quadratic-gradient}
		\nabla\bigl(x^{\transp}Ax\bigr)=\bigl(A+A^{\transp}\bigr)x .
	\end{equation}
	For $n=2$ this is the matrix
	$\begin{psmallmatrix}2a_{11}&a_{12}+a_{21}\\a_{12}+a_{21}&2a_{22}\end{psmallmatrix}$
	applied to $x$, as in Exercise~3.3 and Example~11.2 of
	\textcite{fukuda2026notes} \autocite[slide 49]{fukuda2026calculus}. When $A$ is
	symmetric \eqref{eq:quadratic-gradient} reduces to the familiar $2Ax$; when it is
	not, only the symmetric part $\tfrac12(A+A^{\transp})$ of $A$ affects the
	quadratic form, which is \autoref{rem:symmetric-part}. The identities that
	recur in least squares, in the second-order conditions of
	\autoref{ch:optimization} and in the Kalman recursion of
	\autoref{subsec:prob-kalman} are collected in \autoref{app:matrix-calculus}.
\end{eg}

\section{Differentiability}\label{sec:cn-differentiability}

\begin{definition}[Differentiability and the total derivative]\label{def:total-derivative}
	Let $A\subseteq\R^n$ be open and $f\colon A\to\R$. The function $f$ is
	\dfnas{differentiable at}{differentiability!in several variables} $x_0\in A$ if
	there is a vector $\nabla f(x_0)\in\R^n$ such that
	\begin{equation}\label{eq:frechet}
		\beta_{x_0}(u)\eqdef
		\frac{f(x_0+u)-f(x_0)-\ip{\nabla f(x_0)}{u}}{\norm{u}}
		\;\longrightarrow\;0
		\qquad\text{as }u\to0,\ u\neq0 .
	\end{equation}
	The associated linear functional $Df(x_0)\colon\R^n\to\R$,
	$Df(x_0)u\eqdef\ip{\nabla f(x_0)}{u}$, is the \dfn{total derivative} of $f$ at
	$x_0$. The function is differentiable on $A$ if it is differentiable at every
	point of $A$.
	\nidx{Df}{$Df(x_0)$, total derivative}
\end{definition}

This is Definition~11.1 of \textcite[p.~369]{fukuda2026notes}. Condition
\eqref{eq:frechet} says that $f(x_0)+\ip{\nabla f(x_0)}{u}$ approximates
$f(x_0+u)$ with an error small relative to $\norm{u}$; geometrically, the graph
has a tangent hyperplane at $x_0$ \autocite[slide 50]{fukuda2026calculus}. The
existence of the vector $\nabla f(x_0)$ is part of the requirement, and by
\autoref{thm:riesz-rn} any linear functional on $\R^n$ has such a
representation, so nothing is lost by stating the definition with an inner
product rather than an abstract linear map.

\begin{proposition}[Consistency with the one-variable definition]\label{prop:consistency-1d}
	For $n=1$, \autoref{def:total-derivative} holds at $x_0$ with
	$\nabla f(x_0)=L$ if and only if \autoref{def:derivative} holds with
	$f'(x_0)=L$.
\end{proposition}

\begin{proof}
	For $n=1$, \eqref{eq:frechet} reads
	$\bigl(f(x_0+u)-f(x_0)-Lu\bigr)/\abs{u}\to0$, which since
	$\abs{\abs{u}}=\abs{u}$ is equivalent to
	$\bigl(f(x_0+u)-f(x_0)\bigr)/u\to L$. This is
	\autoref{prop:diff-continuous}, and hence \autoref{def:derivative}.
\end{proof}

\begin{proposition}[Differentiability implies continuity and the existence of partials]
	\label{prop:differentiable-implies}
	If $f$ is differentiable at $x_0$, then $f$ is continuous at $x_0$, every
	partial derivative of $f$ exists at $x_0$, and the vector $\nabla f(x_0)$ in
	\eqref{eq:frechet} is the gradient of \autoref{def:partial}. Moreover the
	directional derivative in any direction $v\in\R^n$ exists and equals
	$\ip{\nabla f(x_0)}{v}$.
\end{proposition}

\begin{proof}
	Continuity: from \eqref{eq:frechet},
	$f(x_0+u)-f(x_0)=\ip{\nabla f(x_0)}{u}+\norm{u}\beta_{x_0}(u)$, and both terms
	tend to $0$ as $u\to0$, the first by Cauchy--Schwarz
	(\autoref{lem:cauchy-schwarz}) and the second because
	$\beta_{x_0}(u)\to0$.

	Directional derivatives: put $u=tv$ with $t\neq0$ and $v\neq0$. Then
	\[
		\frac{f(x_0+tv)-f(x_0)}{t}
		=\ip{\nabla f(x_0)}{v}+\frac{\abs{t}\norm{v}}{t}\,\beta_{x_0}(tv)
		\;\longrightarrow\;\ip{\nabla f(x_0)}{v},
	\]
	because $\abs{\abs{t}/t}=1$ and $\beta_{x_0}(tv)\to0$. Taking $v=e_i$ gives the
	$i$th partial derivative and identifies the $i$th coordinate of
	$\nabla f(x_0)$.
\end{proof}

The converse fails in both of the ways that matter, and the two counterexamples
are worth stating because each corresponds to a real modelling error.

\begin{eg}[Partials without continuity]\label{eg:partials-no-continuity}
	Let $f\colon\R^2\to\R$ be $f(x_1,x_2)=0$ if $x_1x_2=0$ and $f(x_1,x_2)=1$
	otherwise. Along each axis $f$ is identically zero, so
	$\partial f/\partial x_1(0,0)=\partial f/\partial x_2(0,0)=0$. Yet
	$f(\varepsilon,\varepsilon)=1$ for every $\varepsilon\neq0$, so $f$ is not
	continuous at the origin. In one variable this cannot happen
	(\autoref{prop:diff-continuous}); in several, the partial derivatives see only
	$n$ lines through the point and are blind to everything else
	\autocite[Rem.~11.4, p.~375]{fukuda2026notes}.
\end{eg}

\begin{eg}[Continuity and partials without differentiability]\label{eg:no-differentiability}
	Let
	\[
		f(x,y)=
		\begin{cases}
			\dfrac{xy}{\sqrt{x^2+y^2}}, & (x,y)\neq(0,0), \\[1.5ex]
			0,                          & (x,y)=(0,0).
		\end{cases}
	\]
	Since $\abs{xy}\leq\tfrac12(x^2+y^2)$, we have
	$\abs{f(x,y)}\leq\tfrac12\sqrt{x^2+y^2}\to0$, so $f$ is continuous at the
	origin. Both partial derivatives at the origin vanish, since $f$ is identically
	zero on each axis. If $f$ were differentiable at the origin, then by
	\autoref{prop:differentiable-implies} the gradient would be $0$ and
	\eqref{eq:frechet} would require
	\[
		\beta_0(x,y)=\frac{f(x,y)}{\sqrt{x^2+y^2}}=\frac{xy}{x^2+y^2}\to0 .
	\]
	Along $(x,y)=(h,h)$ this ratio equals $\tfrac12$ for every $h\neq0$, so the
	limit is not $0$ and $f$ is not differentiable at the origin
	\autocite[Rem.~11.3, p.~374]{fukuda2026notes}. The obstruction is directional:
	$f$ has a directional derivative in every direction but they do not fit together
	into a linear functional.
\end{eg}

\begin{proposition}[Continuous partials suffice]\label{prop:c1-differentiable}
	Let $A\subseteq\R^n$ be open and let $f\colon A\to\R$ have partial derivatives
	with respect to every variable at every point of $A$, all of them continuous at
	$x_0\in A$. Then $f$ is differentiable at $x_0$.
\end{proposition}

\begin{proof}
	Write $u=(u_1,\dots,u_n)$ and interpolate one coordinate at a time. Set
	$v^0\eqdef x_0$ and $v^k\eqdef x_0+\sum_{i\leq k}u_ie_i$, so $v^n=x_0+u$. Then
	\[
		f(x_0+u)-f(x_0)=\sum_{k=1}^{n}\bigl(f(v^k)-f(v^{k-1})\bigr),
	\]
	and each summand differs only in the $k$th coordinate. Applying the
	one-variable mean value theorem (\autoref{thm:mvt}) to that section gives
	$\xi^k$ on the segment from $v^{k-1}$ to $v^k$ with
	$f(v^k)-f(v^{k-1})=\frac{\partial f}{\partial x_k}(\xi^k)\,u_k$. Hence
	\[
		f(x_0+u)-f(x_0)-\ip{\nabla f(x_0)}{u}
		=\sum_{k=1}^{n}\left(\frac{\partial f}{\partial x_k}(\xi^k)
		-\frac{\partial f}{\partial x_k}(x_0)\right)u_k .
	\]
	Since $\abs{u_k}\leq\norm{u}$ and $\xi^k\to x_0$ as $u\to0$, continuity of the
	partials makes each bracket tend to $0$, so the whole expression divided by
	$\norm{u}$ tends to $0$.
\end{proof}

This is Proposition~11.4 of \textcite[p.~375]{fukuda2026notes}. It supplies the
practical criterion: a function given by an explicit formula built from
elementary operations has continuous partials wherever those operations are
defined, hence is differentiable there, and the pathologies of
\autoref{eg:partials-no-continuity} and \autoref{eg:no-differentiability} arise
only at points where the formula changes.

\begin{definition}[Continuous differentiability]\label{def:c1}
	$f$ is \dfnas{continuously differentiable}{continuously differentiable} on an
	open $A$, written $f\in C^1(A)$, if all partial derivatives exist and are
	continuous on $A$. Inductively, $f\in C^k(A)$ if all partial derivatives of
	order $k$ exist and are continuous.
	\nidx{Ck}{$C^k$, $k$ times continuously differentiable}
\end{definition}

The implications established so far, none of which reverses, are
\[
	f\in C^1\ \Longrightarrow\ f\text{ differentiable}\ \Longrightarrow\
	\begin{cases}
		f\text{ continuous},                \\
		\text{all partials exist}
	\end{cases}
\]
and neither of the two consequences on the right, alone or together, implies
differentiability.

\begin{remark}[The total differential]\label{rem:total-differential}
	Writing $\diff y$ for the increment of $f$ and $\diff x_i$ for the increments of the
	arguments, \eqref{eq:frechet} is the statement
	\[
		\diff y=\pderiv{f}{x_1}\,\diff x_1+\cdots+
		\pderiv{f}{x_n}\,\diff x_n
	\]
	up to an error small relative to $\norm{\diff x}$
	\autocite[slide 51]{fukuda2026calculus}. The classical notation is convenient
	and hides the content: the identity is not an algebraic manipulation of
	infinitesimals but the assertion that the error term in \eqref{eq:frechet}
	vanishes faster than $\norm{u}$, which by
	\autoref{eg:no-differentiability} is a genuine restriction on $f$ and not a
	consequence of the existence of the partial derivatives appearing in it.
\end{remark}

\subsection{Vector-valued maps and the Jacobian}

\begin{definition}[Jacobian]\label{def:jacobian}
	Let $A\subseteq\R^n$ be open and $f=(f_1,\dots,f_m)\colon A\to\R^m$. Then $f$
	is differentiable at $x_0$ if there is a linear map
	$Df(x_0)\colon\R^n\to\R^m$ with
	$\norm{f(x_0+u)-f(x_0)-Df(x_0)u}/\norm{u}\to0$ as $u\to0$. Equivalently, each
	component $f_k$ is differentiable at $x_0$; and the matrix of $Df(x_0)$ in the
	canonical bases is the \dfn{Jacobian matrix}
	\[
		Df(x_0)=\begin{pmatrix}
			\dfrac{\partial f_1}{\partial x_1}(x_0) & \cdots & \dfrac{\partial f_1}{\partial x_n}(x_0) \\[1.5ex]
			\vdots                                  &        & \vdots                                  \\[0.5ex]
			\dfrac{\partial f_m}{\partial x_1}(x_0) & \cdots & \dfrac{\partial f_m}{\partial x_n}(x_0)
		\end{pmatrix}\in M_{m\times n}(\R).
	\]
	\nidx{Jacobian}{$Df$, Jacobian matrix}
\end{definition}

The row-wise reading of a matrix flagged in \autoref{thm:matrix-rep} appears
here: row $k$ of $Df(x_0)$ is $\nabla f_k(x_0)^{\transp}$. For $m=1$ the
Jacobian is the \emph{row} vector $\nabla f(x_0)^{\transp}$, which is why
gradients and Jacobians differ by a transpose and why the chain rule below is
stated for matrices rather than vectors.

\begin{theorem}[Chain rule]\label{thm:chain-rule}
	Let $A\subseteq\R^n$ and $B\subseteq\R^m$ be open, $f\colon A\to\R^m$ with
	$f(A)\subseteq B$, and $\varphi\colon B\to\R^p$. If $f$ is differentiable at
	$x_0$ and $\varphi$ is differentiable at $f(x_0)$, then $\varphi\circ f$ is
	differentiable at $x_0$ and
	\begin{equation}\label{eq:chain-rule}
		D(\varphi\circ f)(x_0)=D\varphi\bigl(f(x_0)\bigr)\,Df(x_0),
	\end{equation}
	a product of a $p\times m$ by an $m\times n$ matrix. Componentwise, for
	$p=1$,
	\begin{equation}\label{eq:chain-rule-components}
		\frac{\partial(\varphi\circ f)}{\partial x_i}(x_0)
		=\sum_{k=1}^{m}\frac{\partial\varphi}{\partial y_k}\bigl(f(x_0)\bigr)\,
		\frac{\partial f_k}{\partial x_i}(x_0).
	\end{equation}
\end{theorem}

\begin{proof}
	Write $y_0=f(x_0)$, $S\eqdef Df(x_0)$ and $T\eqdef D\varphi(y_0)$. By
	differentiability,
	$f(x_0+u)=y_0+Su+\norm{u}\varepsilon_1(u)$ with $\varepsilon_1(u)\to0$, and
	$\varphi(y_0+v)=\varphi(y_0)+Tv+\norm{v}\varepsilon_2(v)$ with
	$\varepsilon_2(v)\to0$. Put $v(u)\eqdef Su+\norm{u}\varepsilon_1(u)$; then
	\[
		\varphi\bigl(f(x_0+u)\bigr)-\varphi(y_0)
		=T\bigl(Su\bigr)+\norm{u}\,T\varepsilon_1(u)
		+\norm{v(u)}\,\varepsilon_2\bigl(v(u)\bigr).
	\]
	Divide by $\norm{u}$. The second term tends to $0$ because $T$ is linear on a
	finite-dimensional space, hence continuous by
	\autoref{prop:linear-continuous}(ii), and $\varepsilon_1(u)\to0$. For the
	third, $\norm{v(u)}\leq\bigl(\norm{S}+\norm{\varepsilon_1(u)}\bigr)\norm{u}$
	stays bounded relative to $\norm{u}$, while $v(u)\to0$ forces
	$\varepsilon_2(v(u))\to0$. Hence the quotient tends to $0$ and
	$D(\varphi\circ f)(x_0)=TS$.
\end{proof}

This is Theorem~11.1 of \textcite[p.~380]{fukuda2026notes}. Identity
\eqref{eq:chain-rule-components} is the rule ``sum over all channels through
which $x_i$ affects the value'', and it is the formal content of every
total-derivative computation in comparative statics: when a policy parameter
enters an equilibrium condition both directly and through an endogenous variable,
\eqref{eq:chain-rule-components} is what enumerates the channels.

\begin{proposition}[Gradients are orthogonal to level sets]\label{prop:gradient-level-set}
	Let $f\colon A\to\R$ be differentiable at $x_0$ and let
	$\gamma\colon(-\delta,\delta)\to A$ be a differentiable curve with
	$\gamma(0)=x_0$ lying in the level set $\{x : f(x)=f(x_0)\}$. Then
	$\ip{\nabla f(x_0)}{\gamma'(0)}=0$. Moreover, among unit vectors $v$, the
	directional derivative $\ip{\nabla f(x_0)}{v}$ is maximised at
	$v=\nabla f(x_0)/\norm{\nabla f(x_0)}$ when $\nabla f(x_0)\neq0$.
\end{proposition}

\begin{proof}
	The function $t\mapsto f(\gamma(t))$ is constant, so by
	\autoref{thm:chain-rule} its derivative
	$\ip{\nabla f(x_0)}{\gamma'(0)}$ vanishes. The second claim is the equality case
	of Cauchy--Schwarz (\autoref{lem:cauchy-schwarz}):
	$\ip{\nabla f(x_0)}{v}\leq\norm{\nabla f(x_0)}$ with equality exactly when $v$
	is the stated unit vector.
\end{proof}

The economic reading is that the gradient of a utility function is orthogonal to
the indifference curve and points in the direction of steepest preference
increase, and that the marginal rate of substitution between goods $i$ and $j$ is
the ratio $\bigl(\partial u/\partial x_i\bigr)\big/\bigl(\partial u/\partial
	x_j\bigr)$ precisely because the level set is orthogonal to the gradient. This is
Proposition~11.6 of \textcite[p.~385]{fukuda2026notes}.

\subsection{Homogeneous functions}

\begin{definition}[Homogeneity]\label{def:homogeneous}
	Let $X\subseteq\R^n$ be a cone, that is $\lambda x\in X$ whenever $x\in X$ and
	$\lambda>0$. A function $f\colon X\to\R$ is \dfnas{homogeneous of degree
		$k$}{homogeneous function} if $f(\lambda x)=\lambda^kf(x)$ for all $x\in X$ and
	$\lambda>0$.
\end{definition}

\begin{theorem}[Euler]\label{thm:euler-homogeneous}
	Let $X\subseteq\R^n$ be an open cone and $f\colon X\to\R$ differentiable. Then
	$f$ is homogeneous of degree $k$ if and only if
	\begin{equation}\label{eq:euler}
		\ip{\nabla f(x)}{x}=\sum_{i=1}^{n}\frac{\partial f}{\partial x_i}(x)\,x_i
		=k\,f(x)\qquad\text{for all }x\in X .
	\end{equation}
	Moreover, if $f$ is homogeneous of degree $k$ then each partial derivative
	$\partial f/\partial x_i$ is homogeneous of degree $k-1$.
\end{theorem}

\begin{proof}
	($\Rightarrow$) Differentiate $f(\lambda x)=\lambda^kf(x)$ with respect to
	$\lambda$ using \autoref{thm:chain-rule}:
	$\ip{\nabla f(\lambda x)}{x}=k\lambda^{k-1}f(x)$. Setting $\lambda=1$ gives
	\eqref{eq:euler}. Differentiating instead with respect to $x_i$ gives
	$\lambda\,\partial_if(\lambda x)=\lambda^k\partial_if(x)$, which is homogeneity
	of degree $k-1$ of $\partial_if$.

	($\Leftarrow$) Fix $x$ and set $g(\lambda)\eqdef\lambda^{-k}f(\lambda x)$ for
	$\lambda>0$. By \autoref{thm:chain-rule} and \eqref{eq:euler} applied at
	$\lambda x$,
	\[
		g'(\lambda)=-k\lambda^{-k-1}f(\lambda x)+\lambda^{-k}\ip{\nabla f(\lambda x)}{x}
		=-k\lambda^{-k-1}f(\lambda x)+\lambda^{-k-1}\,k f(\lambda x)=0 ,
	\]
	using $\ip{\nabla f(\lambda x)}{\lambda x}=kf(\lambda x)$. So $g$ is constant
	on $(0,\infty)$ by \autoref{thm:mvt}, and $g(1)=f(x)$ gives
	$f(\lambda x)=\lambda^kf(x)$.
\end{proof}

Homogeneity of degree $1$ of a production function is constant returns to scale,
and \eqref{eq:euler} then reads $Y=F_KK+F_LL$: output is exhausted by payments to
factors at their marginal products. This is Definition~11.3 of
\textcite[p.~385]{fukuda2026notes}, and it is the identity behind the factor
shares computed in \autoref{sec:cn-firm}.

\section{Optimisation Without Constraints}\label{sec:cn-foc}

\begin{theorem}[Fermat, several variables]\label{thm:fermat-n}
	Let $A\subseteq\R^n$, let $x^\ast$ be an interior point of $A$ at which
	$f\colon A\to\R$ is differentiable, and suppose $x^\ast$ is a local maximiser or
	local minimiser of $f$ on $A$. Then $\nabla f(x^\ast)=0$.
\end{theorem}

\begin{proof}
	Fix $i$ and consider the section $g_i(t)\eqdef f(x^\ast+te_i)$, defined for
	$\abs{t}$ small because $x^\ast$ is interior. Then $g_i$ is differentiable at
	$0$ with $g_i'(0)=\partial f/\partial x_i(x^\ast)$ by
	\autoref{prop:differentiable-implies}, and $0$ is a local extremum of $g_i$. By
	\autoref{thm:fermat}, $g_i'(0)=0$. As $i$ was arbitrary,
	$\nabla f(x^\ast)=0$.
\end{proof}

\begin{theorem}[Sufficiency under concavity]\label{thm:foc-concave-n}
	Let $C\subseteq\R^n$ be convex, $f\colon C\to\R$ continuous on $C$,
	differentiable on $\Int C$, and concave. Then $x^\ast\in\Int C$ is a global
	maximiser of $f$ on $C$ if and only if $\nabla f(x^\ast)=0$. The convex case
	with minimisers is obtained by replacing $f$ with $-f$.
\end{theorem}

\begin{proof}
	Necessity is \autoref{thm:fermat-n}. For sufficiency, the multivariate
	analogue of \autoref{thm:concave-tangent}(ii), proved as
	\autoref{thm:concave-gradient}, gives
	$f(x)\leq f(x^\ast)+\ip{\nabla f(x^\ast)}{x-x^\ast}=f(x^\ast)$ for all
	$x\in C$, first for $x\in\Int C$ and then for $x\in C$ by continuity.
\end{proof}

This is the content of \autocite[slide 52]{fukuda2026calculus}.

\begin{eg}[The four canonical surfaces]\label{eg:four-surfaces}
	At the origin, all four functions below have vanishing gradient
	\autocite[slide 53]{fukuda2026calculus}.
	\begin{enumerate}[(i)]
		\item $f(x,y)=x^2+y^2$: a strict global minimum. $\Hess f=2I\succ0$.
		\item $f(x,y)=-x^2-y^2$: a strict global maximum. $\Hess f=-2I\prec0$.
		\item $f(x,y)=x^2-y^2$: a saddle. $\Hess f=\diag(2,-2)$ is indefinite;
		      the origin is a minimum along the $x$-axis and a maximum along the
		      $y$-axis, so it is neither.
		\item $f(x,y)=-x^3$: neither, and the Hessian at the origin is the zero
		      matrix, which is both positive and negative semidefinite. The second-order
		      test is silent, and the third derivative decides.
	\end{enumerate}
	Cases (iii) and (iv) are the two ways in which a stationary point can fail to be
	an extremum, and they are distinguished by whether the Hessian is indefinite
	(the test rules out an extremum) or merely singular (the test is
	inconclusive).
\end{eg}

\section{Second-Order Theory}\label{sec:cn-second}

\begin{definition}[Hessian]\label{def:hessian}
	Let $A\subseteq\R^n$ be open and $f\colon A\to\R$ twice differentiable at
	$x\in A$. The \dfn{Hessian} of $f$ at $x$ is the matrix of second partial
	derivatives
	\[
		\Hess f(x)\eqdef
		\left(\frac{\partial^2f}{\partial x_i\partial x_j}(x)\right)_{i,j=1}^{n}
		\in M_{n\times n}(\R),
	\]
	also written $Hf(x)$ or $\nabla^2f(x)$; it is the Jacobian of the map
	$x\mapsto\nabla f(x)$.
	\nidx{Hessian}{$\Hess f$, Hessian matrix}
\end{definition}

\begin{theorem}[Schwarz]\label{thm:schwarz}
	If $f\in C^2(A)$ then $\Hess f(x)$ is symmetric for every $x\in A$:
	$\partial^2f/\partial x_i\partial x_j=\partial^2f/\partial x_j\partial x_i$.
\end{theorem}

Symmetry is not automatic without the continuity hypothesis; the standard
counterexample is $f(x,y)=xy(x^2-y^2)/(x^2+y^2)$ with $f(0,0)=0$, for which the
two mixed partials at the origin are $-1$ and $+1$. See
\textcite[Thm.~9.41]{rudin1976principles}. In economics symmetry of the Hessian
is not a technicality: it is the source of the Slutsky symmetry conditions and of
the reciprocity relations that make the Hessian of an expenditure function
testable.

\begin{theorem}[Second-order Taylor expansion]\label{thm:taylor-n}
	Let $A\subseteq\R^n$, let $\bar x\in\Int A$, and let $f\colon A\to\R$ be twice
	differentiable at $\bar x$. Then
	\begin{equation}\label{eq:taylor-n}
		f(x)=f(\bar x)+\ip{\nabla f(\bar x)}{x-\bar x}
		+\tfrac12(x-\bar x)^{\transp}\Hess f(\bar x)(x-\bar x)
		+\norm{x-\bar x}^2\alpha_{\bar x}(x-\bar x),
	\end{equation}
	where $\alpha_{\bar x}(x-\bar x)\to0$ as $x\to\bar x$. If $f\in C^2$ on a
	neighbourhood of the segment $[\bar x,x]$, there is
	$\theta\in(0,1)$ with the exact expression
	\[
		f(x)=f(\bar x)+\ip{\nabla f(\bar x)}{x-\bar x}
		+\tfrac12(x-\bar x)^{\transp}\Hess f\bigl(\bar x+\theta(x-\bar x)\bigr)(x-\bar x).
	\]
\end{theorem}

\begin{proof}
	For the exact form, fix $x$ and set $\varphi(t)\eqdef f\bigl(\bar
		x+t(x-\bar x)\bigr)$ for $t\in[0,1]$. By \autoref{thm:chain-rule},
	$\varphi'(t)=\ip{\nabla f(\bar x+t(x-\bar x))}{x-\bar x}$ and, applying the
	chain rule again to the gradient map,
	$\varphi''(t)=(x-\bar x)^{\transp}\Hess f(\bar x+t(x-\bar x))(x-\bar x)$.
	Applying the one-variable Taylor theorem
	(\autoref{thm:taylor-1d} with $n=1$) to $\varphi$ on $[0,1]$ gives
	$\varphi(1)=\varphi(0)+\varphi'(0)+\tfrac12\varphi''(\theta)$ for some
	$\theta\in(0,1)$, which is the claim. The Peano form \eqref{eq:taylor-n} follows
	from twice differentiability at the single point $\bar x$ by the same argument
	applied with \autoref{cor:peano}.
\end{proof}

This is Theorem~11.5 of \textcite[p.~413]{fukuda2026notes}.

\begin{theorem}[Second-order conditions]\label{thm:second-order-conditions}
	Let $A\subseteq\R^n$ be open, $f\in C^2(A)$, and $x^\ast\in A$ with
	$\nabla f(x^\ast)=0$.
	\begin{enumerate}[(i)]
		\item If $\Hess f(x^\ast)\prec0$, then $x^\ast$ is a strict local maximiser.
		\item If $\Hess f(x^\ast)\succ0$, then $x^\ast$ is a strict local minimiser.
		\item If $\Hess f(x^\ast)$ is indefinite, then $x^\ast$ is neither.
		\item If $x^\ast$ is a local maximiser then $\Hess f(x^\ast)\preceq0$; if a
		      local minimiser then $\Hess f(x^\ast)\succeq0$.
	\end{enumerate}
\end{theorem}

\begin{proof}
	(i) Let $H\eqdef\Hess f(x^\ast)\prec0$. The function $v\mapsto v^{\transp}Hv$
	is continuous on the compact unit sphere, so by
	\autoref{cor:weierstrass-continuous} it attains a maximum $-2\mu$ there, and
	$\mu>0$ by negative definiteness. By homogeneity,
	$u^{\transp}Hu\leq-2\mu\norm{u}^2$ for all $u$. From \eqref{eq:taylor-n} with
	$u=x-x^\ast$ and $\nabla f(x^\ast)=0$,
	\[
		f(x)-f(x^\ast)\leq-\mu\norm{u}^2+\norm{u}^2\alpha(u)
		=\norm{u}^2\bigl(\alpha(u)-\mu\bigr),
	\]
	which is negative for $u\neq0$ small enough that $\abs{\alpha(u)}<\mu$. Part
	(ii) applies (i) to $-f$.

	(iii) Indefiniteness supplies $v_+$ with $v_+^{\transp}Hv_+>0$ and $v_-$ with
	$v_-^{\transp}Hv_-<0$. Restricting $f$ to the line $t\mapsto x^\ast+tv_+$ and
	applying \eqref{eq:taylor-n} shows $f(x^\ast+tv_+)>f(x^\ast)$ for small
	$t\neq0$, and symmetrically along $v_-$; so $x^\ast$ is neither a local maximum
	nor a local minimum.

	(iv) If $x^\ast$ is a local maximiser and some $v$ had $v^{\transp}Hv>0$, the
	argument of (iii) along $v$ would contradict maximality.
\end{proof}

The gap between the sufficient conditions (i)--(ii) and the necessary condition
(iv) is exactly the semidefinite case, and it cannot be closed:
$f(x)=x^3$ and $f(x)=x^4$ both have $f''(0)=0$ and differ in their behaviour.
Definiteness is checked by the eigenvalue criterion
(\autoref{prop:definiteness-eigen}) or by Sylvester's criterion
(\autoref{prop:sylvester}), with the warning of
\autoref{rem:semidefinite-minors} about the semidefinite case.

\begin{theorem}[Hessian characterisation of convexity]\label{thm:hessian-convexity}
	Let $A\subseteq\R^n$ be open and convex and $f\in C^2(A)$. Then $f$ is convex
	on $A$ if and only if $\Hess f(x)\succeq0$ for every $x\in A$, and concave if
	and only if $\Hess f(x)\preceq0$ for every $x\in A$. If $\Hess f(x)\succ0$
	everywhere then $f$ is strictly convex; the converse fails, as $f(x)=x^4$
	shows.
\end{theorem}

\begin{proof}
	Convexity of $f$ on the convex set $A$ is equivalent to convexity of
	$\varphi_{x,y}(t)\eqdef f\bigl(x+t(y-x)\bigr)$ on $[0,1]$ for every
	$x,y\in A$, since the two statements are the same inequality. By
	\autoref{thm:concave-tangent} applied to $-\varphi_{x,y}$, that holds if and
	only if $\varphi_{x,y}''\geq0$, and
	$\varphi_{x,y}''(t)=(y-x)^{\transp}\Hess f(x+t(y-x))(y-x)$ as computed in the
	proof of \autoref{thm:taylor-n}. Ranging over $x,y\in A$ makes $y-x$ an
	arbitrary direction and $x+t(y-x)$ an arbitrary point of $A$, which is exactly
	$\Hess f\succeq0$ on $A$.
\end{proof}

This is Theorem~11.6 of \textcite[p.~415]{fukuda2026notes}.

\section{The Implicit and Inverse Function Theorems}\label{sec:cn-ift}

\begin{theorem}[Implicit function theorem]\label{thm:implicit-function}
	Let $\Lambda\subseteq\R^n\times\R^m$ be open and $f\colon\Lambda\to\R^m$
	continuously differentiable. Let $(\bar x,\bar y)\in\Lambda$ satisfy
	$f(\bar x,\bar y)=0$ and suppose the $m\times m$ matrix
	$\nabla_yf(\bar x,\bar y)$ is nonsingular. Then there are an open neighbourhood
	$\Gamma\subseteq\R^n$ of $\bar x$ and a continuously differentiable map
	$e\colon\Gamma\to\R^m$ such that
	\begin{enumerate}[(i)]
		\item $(x,e(x))\in\Lambda$ and $f\bigl(x,e(x)\bigr)=0$ for all $x\in\Gamma$;
		\item $e(\bar x)=\bar y$;
		\item $\nabla_yf\bigl(x,e(x)\bigr)$ is nonsingular for all $x\in\Gamma$; and
		\item for all $x\in\Gamma$,
		      \begin{equation}\label{eq:ift-derivative}
			      De(x)=-\bigl[\nabla_yf\bigl(x,e(x)\bigr)\bigr]^{-1}
			      \nabla_xf\bigl(x,e(x)\bigr).
		      \end{equation}
	\end{enumerate}
	Moreover $e$ is the unique such function on a possibly smaller neighbourhood:
	the solutions of $f(x,y)=0$ near $(\bar x,\bar y)$ are exactly the points
	$(x,e(x))$.
\end{theorem}

This is Theorem~11.7 of \textcite[p.~419]{fukuda2026notes}. The proof constructs
$e$ as the fixed point of a contraction: for each $x$ near $\bar x$ the map
$y\mapsto y-[\nabla_yf(\bar x,\bar y)]^{-1}f(x,y)$ is a contraction on a small
closed ball, by continuity of $Df$, and \autoref{thm:banach-fixed-point} applies
on that complete set, with \autoref{cor:invariant-set} keeping the fixed point
inside the ball; the regularity of $e$ then follows from the regularity of $f$.
It is carried out in full in \autoref{app:ift}, where it is placed because it is
long and purely technical, and because it is exactly the use of the contraction
mapping theorem advertised on \autocite[slide 65]{fukuda2026analysis}; the
classical treatment is \textcite[Thm.~9.28]{rudin1976principles}.

Formula \eqref{eq:ift-derivative} is not an extra: it is obtained by
differentiating the identity $f(x,e(x))=0$ with respect to $x$ using
\autoref{thm:chain-rule},
\[
	\nabla_xf\bigl(x,e(x)\bigr)+\nabla_yf\bigl(x,e(x)\bigr)De(x)=0,
\]
and solving, which is legitimate precisely because (iii) guarantees the inverse
exists. The differentiation is valid only once $e$ is known to be
differentiable, which is why (iv) is a conclusion of the theorem and not a
derivation of it.

\begin{corollary}[Inverse function theorem]\label{thm:inverse-function}
	Let $A\subseteq\R^n$ be open, $g\colon A\to\R^n$ continuously differentiable,
	and $\bar x\in A$ with $Dg(\bar x)$ nonsingular. Then there are open
	neighbourhoods $U\ni\bar x$ and $V\ni g(\bar x)$ such that
	$\restr{g}{U}\colon U\to V$ is a bijection with continuously differentiable
	inverse, and
	\[
		D\bigl(g^{-1}\bigr)(y)=\bigl[Dg\bigl(g^{-1}(y)\bigr)\bigr]^{-1},
		\qquad y\in V .
	\]
\end{corollary}

\begin{proof}
	Apply \autoref{thm:implicit-function} to $f(y,x)\eqdef g(x)-y$ on
	$\R^n\times A$, at the point $(g(\bar x),\bar x)$. Here the roles are exchanged:
	the ``parameter'' is $y$ and the ``unknown'' is $x$, and
	$\nabla_xf=Dg(\bar x)$ is nonsingular by hypothesis. The implicit function
	$e$ is the local inverse, and \eqref{eq:ift-derivative} gives
	$De(y)=-[Dg]^{-1}(-I)=[Dg]^{-1}$.
\end{proof}

For $n=1$ this is \autoref{prop:inverse-derivative}, with $Dg(\bar x)\neq0$ in
place of $f'(x_0)\neq0$. The theorem is local and cannot be made global: the map
$g(x_1,x_2)=(e^{x_1}\cos x_2,e^{x_1}\sin x_2)$ has everywhere nonsingular
Jacobian and is not injective, being $2\pi$-periodic in $x_2$.

\begin{remark}[Comparative statics is the implicit function theorem]
	\label{rem:comparative-statics-ift}
	The standard comparative-statics exercise has this form. An endogenous vector
	$y\in\R^m$ and an exogenous vector $x\in\R^n$ satisfy $m$ equilibrium conditions
	$f(x,y)=0$ --- first-order conditions, market-clearing conditions, or both. At a
	reference equilibrium $(\bar x,\bar y)$ one asks how $y$ responds to $x$. The
	theorem answers three questions at once, and it is worth separating them.

	\emph{Does an equilibrium exist nearby?} Yes, and it varies continuously: this is
	conclusion (i)--(ii).

	\emph{Is it locally unique?} Yes: this is the uniqueness clause, and it is what
	makes ``the'' comparative static well posed. Uniqueness is local only; the
	theorem says nothing about other equilibria far away.

	\emph{What is the response?} It is \eqref{eq:ift-derivative}. The economic
	content is in the two factors. The matrix $\nabla_xf$ collects the direct
	effects of the parameters on the equilibrium conditions; the matrix
	$[\nabla_yf]^{-1}$ propagates them through the simultaneous system. When
	$f=\nabla_y\Pi$ for an objective $\Pi$, the propagating matrix is the inverse
	Hessian, and \autoref{thm:second-order-conditions} makes it negative definite at
	a maximum --- which is the source of the standard sign restrictions on
	comparative statics.

	The condition that can fail is nonsingularity of $\nabla_yf$. At such a point
	the equilibrium is not locally unique, comparative statics are undefined, and
	the model is at a bifurcation.
\end{remark}

\section{Economic Application: The Representative Firm}\label{sec:cn-firm}

\paragraph{Environment.} A firm rents capital $K\geq0$ at rate $r>0$ and hires
labour $L\geq0$ at wage $w>0$, producing $Y=F(K,L)=AK^{\alpha}L^{1-\alpha}$ with
$A>0$ and $\alpha\in(0,1)$. Output is the numeraire. The firm solves
\eqref{eq:rep-firm}.

\begin{proposition}[First-order conditions and factor shares]\label{prop:rep-firm}
	Let $(K^\ast,L^\ast)$ be an interior solution of \eqref{eq:rep-firm}. Then
	\begin{equation}\label{eq:factor-prices}
		r=\frac{\partial Y}{\partial K}=\alpha A\left(\frac{K^\ast}{L^\ast}\right)^{\alpha-1},
		\qquad
		w=\frac{\partial Y}{\partial L}=(1-\alpha)A\left(\frac{K^\ast}{L^\ast}\right)^{\alpha},
	\end{equation}
	and consequently
	\begin{equation}\label{eq:shares}
		\alpha=\frac{rK^\ast}{Y^\ast},\qquad 1-\alpha=\frac{wL^\ast}{Y^\ast},
		\qquad\text{and}\qquad \Pi(K^\ast,L^\ast)=0 .
	\end{equation}
\end{proposition}

\begin{proof}
	$\Pi$ is differentiable on $\R^2_{++}$ and $(K^\ast,L^\ast)$ is interior, so
	\autoref{thm:fermat-n} gives $\nabla\Pi=0$, which is \eqref{eq:factor-prices}
	after computing the partials as in \autoref{eg:cobb-douglas-partials}.
	Multiplying the first equation by $K^\ast$ gives
	$rK^\ast=\alpha A(K^\ast)^{\alpha}(L^\ast)^{1-\alpha}=\alpha Y^\ast$, and
	similarly for labour. Adding, $rK^\ast+wL^\ast=Y^\ast$, so profit is zero.
\end{proof}

\paragraph{Reading the result.} Three points, each with a mathematical source.

The zero-profit conclusion is \autoref{thm:euler-homogeneous} with $k=1$: $F$ is
homogeneous of degree $1$, so $\ip{\nabla F}{(K,L)}=F(K,L)$, and payment of each
factor at its marginal product exhausts output exactly. Nothing about
Cobb--Douglas is used beyond constant returns.

The factor shares in \eqref{eq:shares} are constants, independent of $r$, $w$ and
$A$. This is special to Cobb--Douglas, and it is the reason the function is used:
it is the production function whose elasticity of substitution between capital
and labour equals $1$, and constancy of shares is equivalent to that.

The scale of the firm is indeterminate. Equations \eqref{eq:factor-prices}
involve $K^\ast$ and $L^\ast$ only through their ratio, because the marginal
products are homogeneous of degree $0$ by the second part of
\autoref{thm:euler-homogeneous}. Given $(r,w)$, either no ratio satisfies both
equations, or one does and then every scalar multiple of $(K^\ast,L^\ast)$ is
also optimal. Formally, the Hessian
\[
	\Hess\Pi=\Hess F=
	\alpha(1-\alpha)\frac{Y}{KL}
	\begin{pmatrix}
		-L/K & 1     \\
		1    & -K/L
	\end{pmatrix}
\]
has determinant $0$ and trace $-\alpha(1-\alpha)\frac{Y}{KL}(L/K+K/L)<0$, so it
is negative \emph{semi}definite with a one-dimensional null space spanned by
$(K,L)$ --- the direction of proportional scaling. The maximum is therefore not
strict, and \autoref{thm:second-order-conditions}(i) does not apply. This is not a
defect of the mathematics but the correct statement that a constant-returns firm
in a competitive market has an indeterminate size: only the aggregate is pinned
down, by demand. This is the optional exercise of
\autocite[slides 43, 54]{fukuda2026calculus}.

\section{Economic Application: Least Squares by Calculus}\label{sec:cn-ols}

\paragraph{Problem.} Given data $(x_i,y_i)_{i=1}^n$ with at least two distinct
$x_i$, minimise the sum of squared residuals \eqref{eq:ols-problem}.

\begin{proposition}\label{prop:ols-calculus}
	Let $S(\beta_1,\beta_2)\eqdef\sum_{i=1}^n(y_i-\beta_1-\beta_2x_i)^2$. If the
	$x_i$ are not all equal, then $S$ has a unique global minimiser, given by
	\begin{equation}\label{eq:ols-solution}
		\hat\beta_2=\frac{\sum_i(x_i-\bar x)(y_i-\bar y)}{\sum_i(x_i-\bar x)^2},
		\qquad
		\hat\beta_1=\bar y-\hat\beta_2\bar x,
	\end{equation}
	where $\bar x=n^{-1}\sum_ix_i$ and $\bar y=n^{-1}\sum_iy_i$.
\end{proposition}

\begin{proof}
	$S$ is a polynomial, hence $C^\infty$. Its partial derivatives are
	\[
		\frac{\partial S}{\partial\beta_1}=-2\sum_i(y_i-\beta_1-\beta_2x_i),
		\qquad
		\frac{\partial S}{\partial\beta_2}=-2\sum_ix_i(y_i-\beta_1-\beta_2x_i),
	\]
	so $\nabla S=0$ is the pair of \dfnas{normal equations}{normal equations}
	\begin{equation}\label{eq:normal-equations}
		n\beta_1+\Bigl(\sum_ix_i\Bigr)\beta_2=\sum_iy_i,
		\qquad
		\Bigl(\sum_ix_i\Bigr)\beta_1+\Bigl(\sum_ix_i^2\Bigr)\beta_2=\sum_ix_iy_i .
	\end{equation}
	The Hessian is
	\[
		\Hess S=2\begin{pmatrix}n & \sum_ix_i \\ \sum_ix_i & \sum_ix_i^2\end{pmatrix},
	\]
	constant in $\beta$, with $\det\Hess S=4\bigl(n\sum_ix_i^2-(\sum_ix_i)^2\bigr)
		=4n\sum_i(x_i-\bar x)^2>0$ when the $x_i$ are not all equal, and leading entry
	$2n>0$. By \autoref{prop:sylvester}, $\Hess S\succ0$, so $S$ is strictly convex
	by \autoref{thm:hessian-convexity} and the unique stationary point is the global
	minimiser by \autoref{thm:foc-concave-n}. Solving \eqref{eq:normal-equations} by
	Cramer's rule (\autoref{eg:cramer}) yields \eqref{eq:ols-solution}.
\end{proof}

\paragraph{What is and is not established.} The calculus route delivers existence,
uniqueness and a formula, and it makes the role of the data condition explicit:
$\sum_i(x_i-\bar x)^2>0$ is simultaneously the nondegeneracy condition on the
regressors, the invertibility condition for the normal equations, and the
positive definiteness condition for the Hessian. It is the two-variable case of
the no-perfect-multicollinearity condition of \autoref{rem:square-matters}.

What the calculus route does not deliver is the interpretation. Formula
\eqref{eq:ols-solution} says that $\hat\beta_2$ is the ratio of a sample
covariance to a sample variance, and \autoref{eg:correlation-cosine} identified
that ratio as a cosine rescaled by norms. The geometry behind it --- that the
fitted values are the orthogonal projection of $y$ onto the span of the
regressors, that the residual is orthogonal to that span, and that the normal
equations \eqref{eq:normal-equations} \emph{are} the orthogonality conditions ---
is developed in \autoref{ch:projection}, where the same estimator is obtained
without differentiating anything and in a form that survives to infinite
dimensions \autocite[slides 55--64]{fukuda2026calculus}.

\section{Chapter Summary}\label{sec:cn-summary}

Partial derivatives differentiate along the coordinate axes; the total derivative
requires a single linear map to approximate the increment in every direction at
once. In one variable the two notions coincide; in several they do not, and the
gap is genuine: a function can have all partial derivatives at a point and be
discontinuous there (\autoref{eg:partials-no-continuity}), or be continuous with
all partials and still fail to be differentiable
(\autoref{eg:no-differentiability}). Continuity of the partials is sufficient
(\autoref{prop:c1-differentiable}), which is why $C^1$ is the standard working
hypothesis. Differentiability delivers continuity, all directional derivatives,
the chain rule in matrix form \eqref{eq:chain-rule}, orthogonality of the
gradient to level sets, and Euler's identity for homogeneous functions.

At an interior local extremum of a differentiable function the gradient vanishes;
under concavity the condition is also sufficient and global. The second-order
theory rests on the Hessian: symmetric under $C^2$ by
\autoref{thm:schwarz}, appearing as the quadratic term of the Taylor expansion
\eqref{eq:taylor-n}, and determining local behaviour at a stationary point
through definiteness, with the semidefinite case inconclusive. Global convexity
or concavity is exactly semidefiniteness of the Hessian everywhere.

The implicit function theorem converts a system of equilibrium conditions into a
locally unique, continuously differentiable solution mapping with derivative
\eqref{eq:ift-derivative}; it is the theorem underlying essentially all
differentiable comparative statics, and its hypothesis --- nonsingularity of the
Jacobian in the endogenous variables --- is exactly what fails at a bifurcation.
Applied to the constant-returns firm, the first-order conditions determine the
capital--labour ratio, exhaust output through Euler's identity, and leave scale
indeterminate, the last fact readable from the singular Hessian.

\begin{exercise}\label{ex:cn-1}
	Let $f(x,y)=x^3y/(x^4+y^2)$ for $(x,y)\neq(0,0)$ and $f(0,0)=0$. Show that the
	directional derivative of $f$ at the origin exists and equals $0$ in every
	direction, yet $f$ is not continuous at the origin. Locate the failure relative
	to \autoref{prop:differentiable-implies}, and explain why this is not a
	counterexample to it.
\end{exercise}

\begin{exercise}\label{ex:cn-2}
	Let $F$ be $C^2$, homogeneous of degree $1$, and strictly quasiconcave on
	$\R^2_{++}$. Show that $\Hess F(K,L)$ is singular with $(K,L)$ in its kernel,
	and deduce that no interior strict local maximum of
	$F(K,L)-rK-wL$ exists. Then show that the cost-minimisation problem
	$\min\{rK+wL : F(K,L)\geq\bar Y\}$ does have a unique solution, and explain the
	difference.
\end{exercise}

\begin{exercise}\label{ex:cn-3}
	A monopolist chooses $q$ facing inverse demand $P(q,\theta)$ with a demand
	shifter $\theta$, and cost $c(q)$. Write the first-order condition as
	$f(\theta,q)=0$ and use \eqref{eq:ift-derivative} to compute
	$dq^\ast/d\theta$. State the exact sign restriction on
	$\partial^2\Pi/\partial q\,\partial\theta$ that makes the response positive, and
	verify that the second-order condition supplies the sign of the denominator.
	Relate the result to \autoref{rem:comparative-statics-ift}.
\end{exercise}
